\documentclass[11pt]{article}
\usepackage[a4paper,margin=26mm]{geometry}
\usepackage{amsmath,amssymb}
\setlength{\parindent}{0pt}
\setlength{\parskip}{7pt}
\emergencystretch=2em
\title{Same network minima, different saddle spectra\\
\large Proof of conjecture 00000006080}
\author{ziangni-sys}
\date{4 October 2026}
\begin{document}
\maketitle
The conjecture asks for two networks with identical minimum positions
and values but different saddle spectra, obtained by an explicit
perturbation. We construct two genuine feedforward networks sharing
both square-activation hidden layers. Changing one output weight
preserves every local and global minimum but changes the Hessian
spectrum of their common saddle.

\textbf{Actual network architecture.}
The activation is $s(u)=u^2$ at each hidden neuron.
For input $(x,y)\in\mathbb R^2$, the first affine layer and activation give
\[
 h_1=\bigl(x^2,(y+1)^2,(y-1)^2\bigr).
\]
The second affine layer uses weights
\[
 W_2=\begin{pmatrix}1&0&0\\0&1/4&-1/4\end{pmatrix},
 \qquad b_2=(-1,0),
\]
so its preactivation is
$\bigl(x^2-1,((y+1)^2-(y-1)^2)/4\bigr)=(x^2-1,y)$.
After squaring, the hidden output is
$\bigl((x^2-1)^2,y^2\bigr)$.
The final linear output weights $(c,1)$ yield
\[
 L_c(x,y)=c(x^2-1)^2+y^2.
\]
The first layer has weights with rows $(1,0),(0,1),(0,1)$
and bias $(0,1,-1)$. Thus every affine layer, activation and output
weight is specified. Our two networks are $L_1$ and $L_2$.
The source fixes no activation class; we use square activation.

\textbf{Every minimum position and value.}
For every $c>0$, $L_c\geq0$, with equality exactly at
\[
 \mathcal P=\{(-1,0),(1,0)\}.
\]
Indeed both summands are nonnegative; their sum vanishes exactly
when $x^2=1$ and $y=0$. Therefore these are precisely all global
minimum positions, each with value zero.

To cover all local minima as well, the actual derivative is
\[
 DL_c(x,y)[u,v]=4cx(x^2-1)u+2yv.
\]
Every local minimum is critical, hence $y=0$ and
$x\in\{-1,0,1\}$.
The origin is not a local minimum: for arbitrarily small
$0<t<1$,
\[
 L_c(t,0)-L_c(0,0)=ct^2(t^2-2)<0.
\]
The remaining two critical points are global minima.
Consequently $L_1$ and $L_2$ have exactly the same local minima,
global minima and minimum values.

\newpage
\textbf{Actual saddle Hessians and full spectra.}
The second derivatives give the actual Hessian at every point:
\[
 H_c(x,y)=
 \begin{pmatrix}4c(3x^2-1)&0\\0&2\end{pmatrix}.
\]
The common critical point $(0,0)$ has
\[
 H_c(0,0)=\begin{pmatrix}-4c&0\\0&2\end{pmatrix}.
\]
For $c>0$, one eigenvalue is negative and the other positive,
so the origin is a strict saddle. Directly, small $x$-axis
perturbations decrease the loss, while nonzero $y$-axis
perturbations increase it.

The characteristic determinant is
\[
 \det(zI-H_c(0,0))=(z+4c)(z-2).
\]
Since a complex square matrix is invertible exactly when its
determinant is nonzero, this characterizes the entire actual
spectrum:
\[
 \sigma(H_1(0,0))=\{-4,2\},\qquad
 \sigma(H_2(0,0))=\{-8,2\}.
\]
These spectra differ. The origin is the only saddle critical
point: the other two critical points are the global minima.

\textbf{Explicit perturbation.}
Only the first output weight changes, from one to two:
\[
 L_2(x,y)-L_1(x,y)=(x^2-1)^2.
\]
More generally the explicit output-weight path $c=1+\theta$,
$0\leq\theta\leq1$, leaves every minimum position and value
unchanged, while the saddle's negative eigenvalue moves
from $-4$ to $-8$. This realizes the requested separation.

\textbf{Formal correspondence.}
Lean defines actual dense affine layers with weight matrices and
biases, square activation and a feedforward network structure.
\texttt{network\_formula} derives the displayed polynomial from
the actual layer evaluation, rather than declaring it a network.

The derivative theorem proves the full \texttt{HasFDerivAt}
formula. The local-minimum theorem and
global-minimum theorem prove the complete minimum-position
sets for every $c>0$, and \texttt{minimum\_value} proves value zero.
The local-minimum exclusion uses an actual arbitrarily close
descent point in each neighborhood.

\texttt{Hessian} is defined by actual second partial derivatives
of its function argument. The Hessian formula computes them,
and the spectrum theorem proves the full Mathlib spectrum
using invertibility and determinants.
\texttt{StrictSaddle} requires a zero Fr\'echet derivative and
negative and positive points in the actual Hessian spectrum.
The final \texttt{counterexample} existentially supplies the
two actual networks, their identical minima and values,
their strict saddles, unequal spectra and explicit perturbation.

\textbf{Reproduction.}
Run \texttt{lake build} in \texttt{lean/} using Lean 4.19.0 and
public Mathlib revision
\begin{center}\small\texttt{c44e0c8ee63ca166450922a373c7409c5d26b00b}\end{center}
The project prints axiom audits.
\end{document}
